\clearpage\begingroup\NoSourceHead\SourceLeaf{534}
\renewcommand{\thefootnote}{\arabic{footnote})}\setcounter{footnote}{0}
\newcommand{\BolzaEqualAboveGreater}{\mathrel{\vcenter{\offinterlineskip\hbox{$\scriptstyle=$}\kern-.15ex\hbox{$>$}}}}
\ArithmeticHead{77}{R}{EIN EXISTENZTHEOREM DER POTENTIALTHEORIE.}
\begin{center}\bfseries III. Teil: Die Arbeit »Allgemeine Lehrsätze in Beziehung auf die im verkehrten\\
Verhältnisse des Quadrats der Entfernung wirkenden Anziehungs-\\
und Abstossungskräfte«\footnotemark.\end{center}
\footnotetext{\textsc{Gauss}' Werke, Bd. V, S. 197.}
\noindent\hspace*{\parindent}1. In der in der Überschrift genannten, im Jahr 1840 erschienenen Arbeit beweist \textsc{Gauss} in Art. 29—34 ein fundamentales \textls[100]{Existenztheorem der
Potentialtheorie} nach einer eigenartigen Methode, die für die Variationsrechnung nach verschiedenen Richtungen hin von Bedeutung ist. Es handelt
sich um den Satz, dass die Dichtigkeit $m$ einer stetigen Massenbelegung einer
gegebenen Fläche $s$ stets und nur auf eine Weise so gewählt werden kann,
dass das Potential
\[
V=\int\frac{m\,ds}{r}
\]
derselben auf der Fläche einer vorgegebenen stetigen Funktion $U$ der Koordinaten eines Punktes der Fläche gleich wird, wobei $r$ den Abstand eines
Punktes des Flächenelementes $ds$ von dem angezogenen Punkt bedeutet und
das Integral über die ganze Fläche zu erstrecken ist.
\par\noindent\hspace*{\parindent}Zum Beweis dieses Satzes geht \textsc{Gauss} von dem folgenden Variationsproblem aus: Unter allen stetigen und gleichartigen Verteilungen (d. h. solchen,
für welche die Dichtigkeit $m$ ihr Vorzeichen auf der Fläche nicht wechselt)
einer gegebenen positiven Masse $M$, für die also
\begin{equation*}\tag{1}
m\BolzaEqualAboveGreater0,\quad\int m\,ds=M,
\end{equation*}
diejenige zu bestimmen, für die das Integral
\[
\Omega=\int(V-2U)m\,ds,
\]
erstreckt über die ganze Fläche $s$, seinen kleinsten Wert annimmt.
\par\noindent\hspace*{\parindent}2. Es wird nun zunächst in Art. 30, 31 gezeigt, dass infolge der Voraussetzung $m\BolzaEqualAboveGreater0$ der Wert des Integrals $\Omega$ für alle zulässigen Verteilungen eine
positive untere Grenze besitzt. Daraus wird dann als selbstverständlich geschlossen, dass das Integral für mindestens eine dieser Verteilungen einen
Minimumswert annehmen muss. Sodann werden mit Hilfe der Variations\mbox{-}
\clearpage\SourceLeaf{535}
\ArithmeticHead{78}{L}{BOLZA, GAUSS UND DIE VARIATIONSRECHNUNG.}
\noindent rechnung die Eigenschaften einer das Minimum liefernden Verteilung abgeleitet, und daraus wird schliesslich die Existenz einer diese Eigenschaften besitzenden Verteilung geschlossen.
\par\noindent\hspace*{\parindent}Wir finden also bei \textsc{Gauss} schon im Jahr 1840 dieselbe Schlussweise
ausgebildet, die später von W. \textsc{Thomson}\footnote{\textit{Sur une équation aux différences partielles qui se présente dans plusieurs questions de physique
mathématique}, Journal de mathématique, Bd. 12 (1847), S. 496.} (1847), \textsc{Dirichlet}\footnote{In Vorlesungen; \textsc{Riemann} sagt darüber im Jahre 1857: »\textsc{Dirichlet} pflegt dieses Prinzip seit einer
Reihe von Jahren in seinen Vorlesungen zu geben«, \textsc{Riemanns} Werke, 2. Aufl., 1892, S. 97.} und \textsc{Riemann}\footnote{\textsc{Riemanns} Werke, 1892, S. 30 und 97 ff.}
(1851) zum Beweis verwandter Existenztheoreme benutzt worden ist, und die
von \textsc{Riemann} als »\textsc{Dirichlet}sches Prinzip«, von den Engländern als »\textsc{Thomson}sches Prinzip« bezeichnet worden ist, aber richtiger den Namen »\textsc{Gauss}sches
Prinzip« führen würde, wie denn auch schon \textsc{Riemann}\footnotemark[3] vermutet hat, dass
\textsc{Dirichlet} zu seinem Prinzip »durch einen ähnlichen Gedanken von \textsc{Gauss}
veranlasst worden ist«. Bekanntlich ist der Schluss von der Existenz einer
unteren Grenze auf die Existenz eines Minimums nicht statthaft, und zwar ist
ja merkwürdiger Weise \textsc{Gauss} selbst der erste gewesen, der bei einer andern
Gelegenheit Bedenken gegen diese Schlussweise ausgesprochen hat, nämlich
in seiner Dissertation\begingroup\renewcommand{\thefootnote}{4a)}\footnote{\textsc{Gauss}' Werke, Bd. III, S. 10. Die betreffende Stelle lautet: »Ex suppositione, $X$ obtinere posse
valorem $S$ neque vero valorem $U$, nondum sequitur, inter $S$ et $U$ necessario valorem $T$ jacere, quem $X$
attingere sed non superare possit. Superest adhuc alius casus: scilicet fieri posset, ut inter $S$ et $U$ limes
situs sit, ad quem accedere quidem quam prope velis possit $X$, ipsum vero nihilominus numquam attingere«.
Mit demselben Einwand beschäftigt sich auch der letzte Absatz von Art. 24 der Dissertation (a. a. O. S. 30)
und ein Brief an \textsc{Schumacher} vom 20. Juni 1840 (Werke, Bd. X 1, S. 108); vgl. die Anmerkung von L. \textsc{Schlesinger}, ebenda, S. 110 und A. \textsc{Fränkel}, \textit{Zahlbegriff und Algebra bei Gauss} (1920), Materialien für eine wiss.
Biogr. von \textsc{Gauss}, Heft VIII, S. 12.}\endgroup\setcounter{footnote}{4}\ bei der Kritik des \textsc{D'Alembert}schen Existenzbeweises
für die Wurzeln einer algebraischen Gleichung. \textsc{Weierstrass}\footnote{K. \textsc{Weierstrass}, »\textit{Ueber das sogenannte Dirichletsche Prinzip}«, Werke, Bd. II, S. 49. Vgl. zur
Geschichte des \textsc{Dirichlet}schen Prinzips auch \textit{Encyklopädie der Mathematischen Wissenschaften}, II A 7 b,
Nr. 17, 23—25.} hat später
(1870) diese Lücke im \textsc{Dirichlet}schen Prinzip besonders nachdrücklich hervorgehoben und durch ein anschauliches Beispiel erläutert. Damit wird dem \textsc{Gauss}schen und den übrigen angeführten Existenzbeweisen die Grundlage entzogen.
\par\noindent\hspace*{\parindent}Ebenso bekannt ist, dass in neuerer Zeit \textsc{Hilbert}\footnote{D. \textsc{Hilbert}, \textit{Über das Dirichletsche Prinzip}, Jahresbericht der Deutschen Mathematikervereinigung,
Bd. 8 (1899), S. 184.} (1899) Methoden
\clearpage\SourceLeaf{536}
\ArithmeticHead{79}{R}{EIN EXISTENZTHEOREM DER POTENTIALTHEORIE.}
\noindent angegeben hat, diese Lücke wieder auszufüllen, indem er gezeigt hat, wie
unter gewissen Voraussetzungen die Existenz einer Lösung eines Variationsproblems a priori bewiesen werden kann.
\clearpage\endgroup
